% ===================================================================
%  TABELL Nr. 3 — Kandheet a Jugend.
%  Traduction 2026 d'apres texto/io, controlee sur texto/fr, texto/de
%  et texto/nl. Consigne : voir texto/lb/10-tabell-01.tex.
%
%  LES DEUX INVERSIONS DU TABLEAU 3 LITUANIEN NE SE PRODUISENT PAS
%  ICI, ET C'EST LE MEME ALINEA QUI LE MONTRE. Le lituanien devait
%  ecrire « varpinės smailę », le clocher devant sa fleche, parce
%  qu'il antepose son genitif ; le luxembourgeois le POSTPOSE avec
%  « vun » — « de Spëtz vum Klackentuerm » — exactement comme l'ido
%  dit « la flecho dil klosheyo ». Meme phrase, meme ordre, aucun
%  tour a chercher.
%
%  IL A POURTANT L'AUTRE CONSTRUCTION SOUS LA MAIN, et c'est ce qui
%  rend le choix interessant. « dem Klackentuerm säi Spëtz » est du
%  luxembourgeois irreprochable, et c'est meme le tour le plus
%  courant a l'oral ; il aurait coute une inversion a chaque page. On
%  prend donc « vun » partout, non pour flatter le releve, mais
%  parce que c'est le tour de l'ecrit, et que le livret est un
%  imprime.
%
%  ET LES MOIS SONT ALLEMANDS QUAND LES JOURS SONT MIXTES. « Mäerz,
%  Abrëll, Mee, Juni, Juli, August » suivent l'allemand ; mais la
%  semaine donne « Méindeg, Dënschdeg, Mëttwoch, Donneschdeg,
%  Freideg, Samschdeg, Sonndeg », ou « Dënschdeg » et « Samschdeg »
%  ne sont ni l'allemand ni le neerlandais. Le calendrier vient du
%  livre, la semaine vient de la bouche — c'est le critere du lieu
%  d'apprentissage, applique a deux listes voisines.
% ===================================================================

%%K t03-apar-1 apar
\VUsaut{3.0mm}
\VUtitre{15.0pt}{60}{TABELL Nr. 3}
\VUsaut{1.6mm}
\VUtitre{11.4pt}{0}{Kandheet a Jugend.}
\VUsaut{3.4mm}

%%K t03-apar-2 apar
(D'3. an d'4. Tabell sinn esou zesummegesat, datt si a véier \VUgras{Familljezeenen}
déi wesentlech Begrëffer iwwer d'\VUgras{Wieder,} d'\VUgras{Alter,} d'\VUgras{Famill,}
d'\VUgras{Kleedung} an de \VUgras{Stand} weisen.)

\VUsaut{4.0mm}
%%K t03-tit-1 sub
\VUcentre{12.0pt}{0}{\textit{Éischt Zeen.}}
\VUsaut{3.0mm}

%%K t03-tit-2 sub
\VUpk{10.2pt}{40}{Daf am Duerf.}
\VUsaut{2.4mm}

%%K t03-tit-3 sub
\VUpk{10.2pt}{40}{D'Fréijoer.}
\VUsaut{3.0mm}

%%K t03-c1-01-1 p
1. --- Mir sinn am \VUgras{Fréijoer,} an de Méint \VUgras{Mäerz}, \VUgras{Abrëll}
oder \VUgras{Mee}, an den Deeg vun der \VUgras{Woch}, \VUgras{Méindeg},
\VUgras{Dënschdeg} oder \VUgras{Mëttwoch}. Et ass \VUgras{zéng Auer} moies.

%%K t03-c1-02-1 p
2. --- D'\VUgras{Wieder} ass schéin, d'\VUgras{Loft} reng. D'Sonn schéngt. E puer
\VUgras{Wollécker} \textsuperscript{(2)} klammen am bloen \VUgras{Himmel}
\textsuperscript{(1)} op.

%%K t03-c1-03-1 p
3. --- D'\VUgras{Beem} hu kaum \VUgras{Blieder.} Déi schlank \VUgras{Pappelen}
\textsuperscript{(3)} an de grousse \VUgras{Platan} \textsuperscript{(17)} hu bis elo
nëmme Kneppercher.

%%K t03-c1-04-1 p
4. --- D'\VUgras{Schwalwen} \textsuperscript{(4)} kommen zréck a fléie mat freedegem
Gepiips ronderëm de \VUgras{Spëtz} \textsuperscript{(7)} vum \VUgras{Klackentuerm}
\textsuperscript{(5)}, deen d'\VUgras{Kierch} \textsuperscript{(6)} vum Duerf kréint.

%%K t03-tit-4 sub
\VUsaut{3.4mm}
\VUpk{10.2pt}{40}{D'Kierch.}
\VUsaut{3.0mm}

%%K t03-c1-05-1 p
5. --- Ganz einfach ass dës Kierch mat hirem groussen \VUgras{Portal}
\textsuperscript{(13)} an hirer déckener \VUgras{Auer} \textsuperscript{(8)}. De
klengen \VUgras{Zeeër} \textsuperscript{(9)} steet op der Zuel X: hie weist
d'\VUgras{Stonnen.} De \VUgras{groussen} \textsuperscript{(10)} steet op XII
(Mëtteg): hie weist d'\VUgras{Minutten.}

%%K t03-c1-06-1 p
6. --- Am Klackentuerm lauden d'\VUgras{Klacken} \textsuperscript{(11)} frou. Uewen
um Daach steet e klengt stengent \VUgras{Kräiz} \textsuperscript{(12)}.

%%K t03-c1-07-1 p
7. --- D'Kierch huet net nëmmen e grousst \VUgras{Portal} \textsuperscript{(13)}, si
huet och eng kleng Säitendier. Duerch dës Dier geet den Här \VUgras{Paschtouer}
\textsuperscript{(15)}, an senger \VUgras{Soutane,} aus der \VUgras{Sakristei}
\textsuperscript{(14)} eraus an an d'\VUgras{Paschtoueschhaus}
\textsuperscript{(19)}.

%%K t03-c1-08-1 p
8. --- Nieft him, do ass d'\VUgras{Mëllechfra} \textsuperscript{(24)} mat hirem
\VUgras{Iesel} \textsuperscript{(23)}. Si kënnt aus der Stad zréck, wou si gutt
Mëllech fir d'Kanner higedroen huet.

%%K t03-tit-5 sub
\VUsaut{4.0mm}
\VUpk{10.2pt}{40}{Den Uebstgaart.}
\VUsaut{3.4mm}

%%K t03-c1-09-1 p
9. --- Den Här Aliboron (den Iesel) ass ganz zefridden mat dësem Moiestour. Hien
otemt mat Genoss d'Loft an, déi vum Doft vun de \VUgras{Bléien} \textsuperscript{(20)}
gebeemt ass, déi d'\VUgras{Uebstbeem} \textsuperscript{(18)} vum nooen
\VUgras{Uebstgaart} bedecken. Ganz rosa a wäiss sinn dës \VUgras{Pëschingsbeem}
\textsuperscript{(18)} an \VUgras{Aprikosebeem} \textsuperscript{(19)}, a si loosse
souguer op eng gutt Recolte hoffen.

%%K t03-c1-10-1 p
10. --- Mëttlerweil ginn déi kleng \VUgras{Beien} \textsuperscript{(22)} aus dem
\VUgras{Beiestack} \textsuperscript{(21)} dohin op Raub, fir summend hire séissen
\VUgras{Hunneg} ze sammelen.

%%K t03-c1-11-1 p
11. --- Mä wat geschitt do? Do lafen d'Kanner vum Duerf erbäi a jagen déi erféierte
\VUgras{Gänsen} \textsuperscript{(25)} ewech. Dat ass de Cortège, dee no der
\VUgras{Daf} vum Jong vum Notaire aus der Kierch erauskënnt.

%%K t03-c1-12-1 p
12. --- De klenge Jhemp, an senger \VUgras{Wéi} \textsuperscript{(26)}, wächt vum
frouen Lauden vun de Klacken op, a seng Schwëster d'Magdalena, déi och de
Daafcortège wëll gesinn, biet hir Mamm, si soll kommen, fir hir d'Hoer ze strähnen
an hir d'Kleeder op der Schwell vum Haus unzedoen.

%%K t03-c1-13-1 p
13. --- D'Mamm ass averstanen. Si leet hire \VUgras{Fichu} \textsuperscript{(28)},
hire \VUgras{Corsage} \textsuperscript{(29)} an hire \VUgras{Schäerz}
\textsuperscript{(30)} un, a si kënnt sech op e klenge Stull virun der Dier setzen.
D'Magdalena huet nëmmen hiren \VUgras{Ënnerrock} \textsuperscript{(31)} an hire
\VUgras{Corset} \textsuperscript{(32)} un.

%%K t03-c1-14-1 p
14. --- Wéi glécklech ass si! Si vergësst hir Léiflingsblummen, hir
\VUgras{Neelercher} \textsuperscript{(34)}, op deenen d'\VUgras{Päiperlecken}
\textsuperscript{(35)} sech setzen, hir \VUgras{Tulpen} \textsuperscript{(36)}, hir
\VUgras{Meiblummen} \textsuperscript{(37)} an \VUgras{Dëppen} \textsuperscript{(38)}
op der \VUgras{Mauer} \textsuperscript{(33)}, an hir \VUgras{Rosen}
\textsuperscript{(39)}, déi esou eng schéin Heck bilden. Si kuckt d'räich Kleedung
vun den Damme vum Cortège un.

%%K t03-c1-15-1 p
15. --- D'Kanner vum Duerf kucken net vill op d'Kleedung. Si stierzen no vir a
stousse sech géigesäiteg, fir \VUgras{Bonbongen} \textsuperscript{(55)} oder
\VUgras{Sue-Stécker} \textsuperscript{(52)} ze kréien. Ee vun hinne kräischt
\textsuperscript{(40)}.

%%K t03-c1-16-1 p
16. --- En anere koum dem \VUgras{Hond} \textsuperscript{(42)} op de Fouss, deen
jäizend fortleeft an d'\VUgras{Hong} \textsuperscript{(43)} a seng
\VUgras{Kichelcher} \textsuperscript{(44)} ewechjage geet.

%%K t03-tit-6 sub
\VUsaut{5.0mm}
\VUpk{10.2pt}{40}{Den Daafcortège.}
\VUsaut{4.0mm}

%%K t03-c1-17-1 p
17. --- Virum Cortège geet d'\VUgras{Kannerfra} \textsuperscript{(45)}. Si huet eng
schéin \VUgras{Kap} \textsuperscript{(46)} mat laange Bänner. Si dréit
d'\VUgras{Neigebuerenes} \textsuperscript{(47)}, dat ganz a \VUgras{Spëtzen}
\textsuperscript{(48)} gekleet ass.

%%K t03-c1-18-1 p
18. --- Hannert der Kannerfra kommen de \VUgras{Peiter} \textsuperscript{(49)} an
d'\VUgras{Giedel} \textsuperscript{(53)}. Si deelen d'Bonbongen an d'Sue-Stécker aus.
De Peiter huet \VUgras{Hänschen} \textsuperscript{(51)} an e \VUgras{Zylinder}
\textsuperscript{(50)}. D'Giedel hält an der Hand e schéine \VUgras{Bouquet}
\textsuperscript{(54)}.

%%K t03-c1-19-1 p
19. --- Hannert hinne gesi mer de \VUgras{Monni} \textsuperscript{(56)} an
d'\VUgras{Tatta} \textsuperscript{(58)} vum Kand. D'Tatta huet en elegante
\VUgras{Hutt} \textsuperscript{(59)}, de Monni e schwaarze \VUgras{Redingote}
\textsuperscript{(57)} an e wäisse Gilet. Duerno kommen de \VUgras{Cousin}
\textsuperscript{(60)} an d'\VUgras{Cousine} \textsuperscript{(61)}. Dës hält bei der
Hand d'\VUgras{Schwëster} \textsuperscript{(62)} vum Neigebuerenen, déi kleng Berta.

%%K t03-c1-20-1 p
20. --- Den anere \VUgras{Brudder} \textsuperscript{(63)} vun der Berta geet hannen.
Hie gëtt d'Hand enger \VUgras{Verwandter} \textsuperscript{(64)}. D'\VUgras{Frënn}
\textsuperscript{(65)} vun der Famill kommen duerno.

%%K t03-tit-7 sub
\VUsaut{4.6mm}
\VUpk{10.2pt}{40}{D'Zuschauer.}
\VUsaut{3.4mm}

%%K t03-c1-21-1 p
21. --- Nieft dem Cortège, do ass e \VUgras{Bettler} \textsuperscript{(66)}, dee sech
op seng \VUgras{Kréck} \textsuperscript{(67)} stäipt. Hie freet no engem Almosen (hie
bettelt).

%%K t03-c1-22-1 p
22. --- Op der anerer Säit ass e klenge ganz \VUgras{héiflechen}
\textsuperscript{(68)} Bouf. Hie gréisst a wënscht de Leit, déi hie kennt, e
\VUgras{gudde Moien}.

%%K t03-c1-23-1 p
23. --- D'Duerfleit si vun doheem erauskomm, fir den Daafcortège ze gesinn. Mir
gesinn e \VUgras{Bauer} \textsuperscript{(69)} mat senger \VUgras{Kattounskap} an
niewendrun eng \VUgras{Baierin} \textsuperscript{(70)}. Duerno eng \VUgras{al Fra}
\textsuperscript{(71)}, déi sech op hire \VUgras{Stack} \textsuperscript{(72)}
stäipt. Schlussendlech de \VUgras{Millermeeschter} \textsuperscript{(73)} mat sengem
groussen \VUgras{Filzhutt} \textsuperscript{(74)} an eng jonk Baierin an
\VUgras{Holzschong} \textsuperscript{(75)}, déi hirem klenge Kand d'Goen bäibréngt.

%%K t03-c1-24-1 p
24. --- Dës Zeen ass ganz ameséiereg.

\VUsaut{4.0mm}
%%K t03-tit-8 sub
\VUcentre{12.0pt}{0}{\textit{Zweet Zeen.}}
\VUsaut{3.0mm}

%%K t03-tit-9 sub
\VUpk{10.2pt}{40}{Ëffentlecht Fest.}
\VUsaut{2.4mm}

%%K t03-tit-10 sub
\VUpk{10.2pt}{40}{Waasserspiller a Regatten.}
\VUsaut{3.0mm}

%%K t03-c2-01-1 p
1. --- De \VUgras{Summer} ass komm. Déi waarm Sonn vum Mount \VUgras{Juni} (Juli oder
\VUgras{August}) dreift déi Jonk un, bueden a rudder ze goen.

%%K t03-c2-02-1 p
2. --- Um rietse Ufer vum Floss zéien d'Rudderer sech aus, fir un de Waasserspiller
an de Regatten deelzehuelen.

%%K t03-c2-03-1 p
3. --- Ee vun hinnen zitt seng \VUgras{Schong} \textsuperscript{(1)} aus; hien
zerbënnt (knäppt) se op. Seng zwee Kameraden si scho prett. Vun elo un hu si nëmmen
nach hiren \VUgras{Tricot} \textsuperscript{(2)} an hir \VUgras{Bodhosen}
\textsuperscript{(3)}. Si schwätzen, während si op hir Frënn waarden. Si schwätzen
iwwer d'Kiermes an d'Fest, dat um anere Ufer ass.

%%K t03-c2-04-1 p
4. --- Um Buedem, nieft hinnen, leien d'\VUgras{Kleeder} vun hire Kameraden: e
\VUgras{Paar} \VUgras{Botten} \textsuperscript{(4)}, \VUgras{Schong}
\textsuperscript{(5)}, \VUgras{Strëmp} \textsuperscript{(6)}, \VUgras{Filzhutt}
\textsuperscript{(7)} a \VUgras{Stréihutt} \textsuperscript{(8)} an eng
\VUgras{Krawatt} \textsuperscript{(9)}.

%%K t03-c2-05-1 p
5. --- Ee vun de jonke Männer huet säi \VUgras{Vestong} \textsuperscript{(10)}
suergfälteg zesummegeluecht. Hie leet e op en Handduch, an dat säin Noper d'zwee
\VUgras{Kleeder} bal aroulen wäert.

%%K t03-c2-06-1 p
6. --- E sechste Bouf ass ze spéit. Hien huet nach seng \VUgras{Kap}
\textsuperscript{(11)} um Kapp. Hien huet weder säin \VUgras{Hiem}
\textsuperscript{(12)}, nach säi \VUgras{Kolli} \textsuperscript{(13)}, nach seng
\VUgras{Manchetten} \textsuperscript{(14)} ausgedoen. Hien hält an der Hand säi
\VUgras{Gilet} \textsuperscript{(17)} an huet nach seng \VUgras{Box}
\textsuperscript{(16)} a seng \VUgras{Bretellen} \textsuperscript{(15)} un. Sécher
wäert hien net prett si fir dee nächste Laf.

%%K t03-c2-07-1 p
7. --- Nieft de jonke Männer féiert e Wee, op deem senge Säite vill
\VUgras{Dëschelen} \textsuperscript{(18)} stinn, bis un d'Ufer vum Baach, wou de
Papp Jhemp tëscht de \VUgras{Rouer} \textsuperscript{(19)} d'\VUgras{Freedefeier}
\textsuperscript{(20)} virbereet, dat owes ofgeschoss gëtt.

%%K t03-tit-11 sub
\VUsaut{4.4mm}
\VUpk{10.2pt}{40}{De Laf.}
\VUsaut{3.4mm}

%%K t03-c2-08-1 p
8. --- Um Baach spadséieren e puer Damme mam \VUgras{Boot} \textsuperscript{(21)},
ënner hire Sonneschiirmer. Si verfollegen de Laf, dee ganz interessant ass. An all
\VUgras{Yol} \textsuperscript{(22)} rudderen véier \VUgras{Rudderer}
\textsuperscript{(24)} mat all hire Kräften, während de \VUgras{Steiermann}
\textsuperscript{(26)} d'\VUgras{Steier} \textsuperscript{(25)} hält an dat
zerbriechlecht Boot féiert. D'\VUgras{Rudderen} \textsuperscript{(23)} schloen am
Takt op d'Waasser, an déi liicht Booter fuere mat schwindelegem Tempo.

%%K t03-c2-09-1 p
9. --- E puer jonk Männer moossen sech, nieft dem Bréckepilier, ëm de Präis vum
\VUgras{Bugspriet} \textsuperscript{(28)}. Ee vun hinne geet virsiichteg op deem
biegsamen a rutschegen Mast, fir dat klengt \VUgras{Fändelchen}
\textsuperscript{(29)} ze paken, dat um Enn befestegt ass. Säin ongléckleche
\VUgras{Konkurrent} \textsuperscript{(30)} ass an d'Waasser gefall. Hie schwëmmt un
d'Ufer.

%%K t03-c2-10-1 p
10. --- Méi al Jongen \VUgras{maachen e Bootstournoi} \textsuperscript{(32)} nieft
dem \VUgras{Kai} \textsuperscript{(33)}. All dës Spiller kuckt e groussen
\VUgras{Publikum} \textsuperscript{(31)}, dee sech op d'\VUgras{Geleeder} vun der
\VUgras{Bréck} \textsuperscript{(27)} stäipt an um \VUgras{Ufer} zesummegedrängt ass.
E puer Zuschauer dréien sech awer ëm, fir d'Spill vum \VUgras{Präismast}
\textsuperscript{(45)} ze verfollegen oder de \VUgras{Cortège vum Buergermeeschter}
ze bewonneren, dee fir d'\VUgras{Verdeelung} vun de Präisser kënnt.

%%K t03-c2-11-1 p
11. --- Fir d'éischt, do sinn e puer \VUgras{Bouwen} \textsuperscript{(35)}, déi de
\VUgras{Museker} \textsuperscript{(36)} virgoen; duerno kommen de
\VUgras{Buergermeeschter} \textsuperscript{(37)}, d'Schäffen an de
\VUgras{Gemengerot} vum Stiedchen. Schlussendlech marschéieren d'\VUgras{Pompjeeën}
\textsuperscript{(38)}.

%%K t03-tit-12 sub
\VUsaut{5.0mm}
\VUpk{10.2pt}{40}{D'Kiermes.}
\VUsaut{3.6mm}

%%K t03-c2-12-1 p
12. --- Op der \VUgras{Foireplaz} \textsuperscript{(39)} sinn vill Leit. Ee kënnt déi
verschidde \VUgras{Baraken} kucken: d'\VUgras{Holzpäerd} \textsuperscript{(40)}, den
\VUgras{Zirkus} \textsuperscript{(41)}, wou d'Virstellung scho ugefaangen huet; den
\VUgras{ëffentleche Bal} \textsuperscript{(42)}, wou déi jonk Männer an déi jonk
Fraen aus der Stad d'Freed vum \VUgras{Danz} genéissen.

%%K t03-c2-13-1 p
13. --- Hannen gesi mer d'\VUgras{Arena} \textsuperscript{(44)}, wou de brave
spuenesche \VUgras{Matador} géint de rosen \VUgras{Stéier} \textsuperscript{(45)}
kämpft, duerno d'\VUgras{Rutschbunn} \textsuperscript{(49)}, d'\VUgras{Schéissbud}
\textsuperscript{(46)}, wou ee sech üübt, \VUgras{Schéissgewierer} ze benotzen an op
d'\VUgras{Ziel} ze schéissen.

%%K t03-c2-14-1 p
14. --- Nieft drun, do ass d'\VUgras{Foiretheater} \textsuperscript{(47)}, wou
d'\VUgras{Comédie} an d'\VUgras{Drama} gespillt ginn. Ganz no drun ass d'Barak vum
\VUgras{Ris} \textsuperscript{(48)} a vum \VUgras{Zwerg.} D'\VUgras{Gréisst} vum
éischten ass zwee Meter fofzéng Zentimeter; deen zweeten ass kaum esou grouss wéi e
Kand.

%%K t03-c2-15-1 p
15. --- Schlussendlech, do ass déi grouss \VUgras{Menagerie} \textsuperscript{(50)}
mat hire \VUgras{wëllen Déieren,} mat hirem \VUgras{Tiger} \textsuperscript{(51)} mat
mächtege \VUgras{Krallen,} mat hirem \VUgras{Léiw} \textsuperscript{(52)} mat
dichter \VUgras{Mähn,} mat hiren zwou \VUgras{Giraffen} \textsuperscript{(53)} a mat
hirem \VUgras{Elefant} \textsuperscript{(54)}, dee säi \VUgras{Rüssel} esou geschéckt
benotzt.

%%K t03-c2-16-1 p
16. --- De \VUgras{Dompteur} \textsuperscript{(55)}, deen dës Déieren dresséiert,
steet virun der Dier vun der Barak.

\VUsaut{6.0mm}
\VUfilet{22mm}
